\documentclass[a4paper, 10pt]{article}

% NALOGA 1
% Vključite pakete za podporo slovenščini,
% prepoznavo vhodnega kodiranja `utf8` in izhodnega kodiranja `T1`.
% Prekopirajte ustrezne tri vrstice iz pregledne datoteke `1-osnove~.tex`.

\usepackage{amsmath}
\usepackage[utf8]{inputenc}
\usepackage[T1]{fontenc}      % izhodna datoteka je kodirana v sistemu T1
\usepackage[slovene]{babel}
\usepackage{amsthm}

\begin{document}

% Pri tej nalogi (in tudi vseh naslednjih) si pomagajte s paleto ukazov!
% Naslov, avtor, datum in povzetek naj bodo izdelani z uporabo ukazov in
% okolja, ki so v ta namen definirani v razredu `article`. Da se bo glava
% dokumenta s temi podatki prikazala, morate uporabiti ukaz `\maketitle`.

\title{Izdelava urnikov}
\author{Nejc Širovnik}
\date{14. \ 10. \ 2019}
\maketitle

{\theoremstyle{plain}
\newtheorem{izrek}{Izrek}
}


{\theoremstyle{definition}
\newtheorem{definicija}{Definicija}
}

\end{document}